\documentclass[12pt,a4paper]{article}

\usepackage[T1]{fontenc}
\usepackage[utf8]{inputenc}
\usepackage[brazil]{babel}
\usepackage{newtxtext,newtxmath}
\usepackage{microtype}
\usepackage{geometry}
\usepackage{setspace}
\usepackage{graphicx}
\usepackage{booktabs}
\usepackage{enumitem}
\usepackage[hidelinks]{hyperref}
\usepackage[round,authoryear]{natbib}

\geometry{margin=2.5cm}
\setstretch{1.15}
\setlength{\parindent}{1.25cm}
\setlength{\parskip}{0pt}

\title{Recuperação aumentada por geração para análise auditável de discursos do Senado Federal}
\author{Fabricio Santana\\
  \small E-mail: \texttt{[PREENCHER]}\\
  \small Órgão/área de atuação: [PREENCHER]}
\date{}

\begin{document}
\maketitle

\begin{abstract}
Este short paper propõe uma solução de recuperação aumentada por geração
(\textit{retrieval-augmented generation}, RAG) para apoiar a consulta auditável
a pronunciamentos do Senado Federal. A proposta combina recuperação de
informação, geração condicionada às evidências recuperadas e apresentação dos
metadados das fontes. Busca-se reduzir o esforço necessário para localizar e
sintetizar manifestações parlamentares sem substituir a análise humana nem
atribuir ao modelo a verificação autônoma de fatos ou de conformidade. O corpus
é composto pelos discursos da 56ª Legislatura, obtidos de fonte pública e
disponibilizados em formato estruturado no Hugging Face. A avaliação proposta
abrange qualidade dos dados, desempenho da recuperação e fidelidade das
respostas às evidências. Espera-se produzir um mecanismo reprodutível de apoio
à transparência e à accountability legislativa, acompanhado de trilhas de
auditoria e salvaguardas contra respostas não fundamentadas.

\noindent\textbf{Palavras-chave:} accountability; discursos parlamentares;
recuperação de informação; RAG; inteligência artificial generativa.
\end{abstract}

\section{Contextualização e descrição do problema}

O Senado Federal disponibiliza dados de pronunciamentos parlamentares em seu
portal de dados abertos \citep{senadoDadosAbertos}. Embora a publicidade dos
registros seja indispensável, o volume textual e a dispersão dos metadados
dificultam a localização e a comparação manual de manifestações relacionadas a
um tema, período ou parlamentar. Buscas estritamente lexicais também podem não
recuperar documentos semanticamente próximos que empreguem vocabulários
diferentes.

Nesse contexto, propõe-se investigar como uma solução RAG auditável pode apoiar
pesquisadores, servidores e instâncias de controle na consulta a pronunciamentos
da 56ª Legislatura. O sistema terá natureza assistiva: deverá recuperar trechos,
sintetizar apenas as evidências disponíveis e apresentar autor, data, partido,
unidade da Federação e endereço da fonte. Não se pretende inferir intenção,
determinar a veracidade de declarações nem automatizar juízos sobre parlamentares.

\section{Fundamentação e técnica selecionada}

A arquitetura RAG associa um mecanismo de recuperação documental a um modelo
gerador, condicionando a resposta ao contexto selecionado
\citep{lewis2020rag}. Em comparação com a consulta manual, a recuperação
semântica amplia as formas de expressão aceitas; em comparação com o uso isolado
de um modelo de linguagem, a recuperação oferece evidências verificáveis e
reduz a dependência do conhecimento paramétrico do modelo.

A escolha é adequada ao problema porque o objeto de análise é essencialmente
textual e já contém metadados úteis à rastreabilidade. A auditabilidade, porém,
depende de controles adicionais: versionamento do corpus, registro dos
parâmetros, citações obrigatórias, possibilidade de abstenção e avaliação
separada da recuperação e da geração.

\section{Metodologia proposta e dados}

O corpus será obtido diretamente do conjunto
\path{fabriciosantana/discursos-senado-legislatura-56}, publicado no Hugging
Face \citep{datasetDiscursos}. A base contém metadados dos pronunciamentos e o
texto integral quando sua coleta foi bem-sucedida. Antes da indexação serão
avaliados esquema, duplicidades, campos ausentes, cobertura textual, datas e
consistência dos identificadores. A revisão também distinguirá falha de coleta
de inexistência de discurso, evitando interpretar ausência de texto como
ausência de atividade parlamentar.

O fluxo proposto compreende: (1) download e registro de versão e hash do
arquivo; (2) validação e normalização dos dados com Python e pandas; (3) seleção
do texto e dos metadados; (4) segmentação dos discursos em trechos com
sobreposição; (5) indexação lexical e vetorial; (6) recuperação híbrida dos
trechos mais relevantes; (7) geração de resposta condicionada às evidências; e
(8) apresentação das citações e armazenamento da trilha de execução.

A avaliação utilizará perguntas definidas previamente e respectivos documentos
relevantes. A etapa de recuperação será mensurada, conforme a disponibilidade
de julgamentos de relevância, por \textit{Recall@k}, precisão em $k$ e
\textit{Mean Reciprocal Rank}. A geração será examinada quanto à correção das
citações, cobertura da pergunta, fidelidade aos trechos e capacidade de
abstenção. Uma amostra será submetida à conferência humana, pois avaliações
automáticas não substituem a inspeção das evidências.

Os dados são públicos, mas o tratamento deverá observar finalidade,
necessidade, integridade e segurança. A solução manterá os textos em seu
contexto documental, respeitará a Lei de Acesso à Informação
\citep{brasil2011lai} e considerará os princípios da Lei Geral de Proteção de
Dados \citep{brasil2018lgpd}. Perfis ou classificações individuais não serão
produzidos sem base metodológica e finalidade legítima.

\section{Resultados esperados e impacto}

Espera-se obter um protótipo capaz de responder a questões temáticas com links
e metadados que permitam ao usuário conferir cada afirmação. Os principais
entregáveis serão o corpus versionado, o índice de recuperação, um conjunto de
perguntas de avaliação, métricas reprodutíveis e registros das evidências usadas
em cada resposta.

A aplicação pode reduzir o tempo de pesquisa e ampliar a cobertura documental
de análises legislativas. Seus resultados devem ser entendidos como subsídio,
e não como prova autônoma. Entre as limitações estão discursos sem texto
integral, erros nos metadados, perda de contexto durante a segmentação, falhas
de recuperação, vieses do corpus e alucinações do modelo gerador. Falsos
negativos podem omitir pronunciamentos relevantes, enquanto falsos positivos
podem associar trechos apenas superficialmente relacionados à pergunta.

\section{Considerações finais}

A proposta articula transparência legislativa, governança de dados e IA
generativa em um fluxo que privilegia a rastreabilidade. Sua viabilidade decorre
da existência de dados públicos estruturados e de ferramentas maduras para
preparação, indexação e avaliação. Como próximos passos, serão construídos o
conjunto de referência, a comparação entre recuperação lexical, vetorial e
híbrida e os critérios de abstenção. A adoção institucional dependerá ainda de
validação com usuários, monitoramento contínuo e definição explícita de
responsabilidades sobre as conclusões produzidas.

\bibliographystyle{plainnat}
\bibliography{referencias}

\end{document}
